\subsection{Cobertura propia y dimensionamiento de atención}
Desde el mes 13 hasta el 56 se reservan un puesto NOC y un puesto SOC continuamente, con escalamiento N2/N3 de ingeniería y reemplazo sin abandonar el resto de los sitios. La mesa bilingüe contará con tres agentes simultáneos de 06:00 a 24:00 en temporada alta (15 de noviembre--31 de marzo) y de 08:00 a 20:00 el resto del año. Fuera de esa franja, NOC recibe inmediatamente emergencias y alertas de navegación vencida por un canal separado 24$\times$7$\times$365, con guardia N2/N3 y prohibición de diferirlas al día siguiente. El dimensionamiento incorpora seis personas NOC, seis SOC y doce de mesa, 24 puestos propios adicionales; los ocho titulares y los refuerzos de ingeniería no cubren esos turnos como una segunda jornada.

El universo de diseño utiliza la proyección del caso: 88 empleados, 170 contratistas, 350 armadores y 590 menores. Para la mesa se reserva una plaza de contacto adulto por menor: 1.198 plazas de contacto de diseño, sin afirmar que existan 590 apoderados distintos. Se supone un 2\,\% de contactos diarios y se concentra conservadoramente ese total en tres horas: $\lambda=1.198(0,02)/3=7,987$ contactos/h; TMA 6 minutos, carga $a=0,799$ Erlang. Son supuestos de oferta, no demanda medida. Erlang C da con tres agentes un 95,42\,\% aproximado antes de 20 segundos; el registro asociado calcula el valor exacto y la sensibilidad a dos y tres veces la demanda. Se validará el pronóstico en marcha blanca mediante llegadas, TMA, abandono y resolución; 80\,\% antes de 20 s, abandono $\leq5$\,\% y resolución inicial $\geq70$\,\% son compromisos, no resultados derivados del modelo. Si el pronóstico exige cuatro agentes, se añadirá ese puesto y su relevo antes de degradar el SLA.\parencite[RT-21.06/14]{bases-transversales}\parencite[RT-21.06; cap.~14]{caso07}

Por mes, la cobertura registra $24D_m$ HH NOC y $24D_m$ SOC, más $3(18D_{a,m}+12D_{b,m})$ de mesa. El pico mensual de 31 días de temporada alta es 744/744/1.674 HH; las capacidades de seis/seis/doce personas a 144 HH son 864/864/1.728. Las diferencias permiten relevo y ausencias, y se contrastan con el cuadrante efectivo; no se imputan los mismos trabajadores a ingeniería o terreno. La jornada y descansos se ajustarán a normativa y calendario, con sustitución por ausencia. Los viajes a boyas dispondrán de embarcación y ventana segura, especialista y suplente; los costos de traslado están incluidos y ninguna salida deja un turno sin cobertura.

La cobertura continua de 13--20 agrega 22.692 HH antes de la fase contractual de Operación; 21--56 conserva 99.378 HH. Son 122.070 HH de turnos en total. Los centros se habilitarán en la sede de Valparaíso antes de 13 y la asistencia de campo se coordinará desde Puerto Montt; su ubicación y disponibilidad se recibirán antes de activar el servicio. Las suplencias cubrirán permisos, vacaciones y enfermedad antes de la ausencia: los 24 puestos son la dotación de rotación mínima y toda ausencia que exceda su capacidad libre se cubrirá con reemplazo propio adicional, sin recargar los turnos. El costo de esa sustitución queda a cargo de Synaptix.\parencite[RT-21.01/03/05/16]{bases-transversales}

La sensibilidad exige cuatro agentes simultáneos con el doble de llegadas y cinco con el triple para mantener al menos 80\,\% antes de 20 segundos. En un mes alto de 31 días, corresponden 2.232 y 2.790 HH de mesa y una rotación mínima por capacidad mensual de dieciséis y veinte personas, respectivamente; el cuadrante deberá confirmar además descansos y ausencias. Esos refuerzos son escenarios de respuesta al pronóstico, separados de la dotación base de doce personas y sin atribuir contratación. Un cuarto agente no basta para el escenario triple. La decisión de aumentar capacidad se tomará con la evidencia de marcha blanca antes de activar la franja de atención y se sostendrá durante el contrato si la demanda lo exige.
